\scanpage{236}\mspage{91}
pour chacune des trois conditions $\mathrm{PS}_g1$, $\mathrm{PS}_g1'$, $\mathrm{PS}2$
séparément, on voit immédiatement que si
elle est vérifiée pour $Y$ et pour les $\,{}_{y\backslash}X$,
elle l'est pour $X$, d'où la proposition.

\uline{Corollaire} 7.3. Les quatre notions envisagées
dans 7.1. sont \deleted{\unsure{filtrantes}}\inserted{stables} par composition.

\diagram{\begin{tikzcd}[ampersand replacement=\&,row sep=small,column sep=small]
X\arrow[d] \& {}_{z\backslash}X\arrow[l]\arrow[d]\\
Y\arrow[d] \& {}_{z\backslash}Y\arrow[l]\arrow[d]\\
Z \& {}_{z\backslash}Z\arrow[l]
\end{tikzcd}}
Démonstration immédiate, en
utilisant la remarque 7.2.

\uline{Remarque} \unsure{7.4.|?.4.} 1) Notons que
la notion de morphisme cofiltrant
\inserted{ou fbt cofiltrant}
\deleted{(et les variantes 7.1)} n'est \uline{pas} stable par
chgt. de base quelconque, ni $=$ par une
\unsure{imm.} ouverte. Exemple : $e\unsure{-|\to}\Delta^1$, appliquant
$e$ en $1$, est cofiltrant, mais sa
fibre en $0$ est vide donc n'est
ni cofiltrante, ni fbt. cofiltrante.
\diagram{\begin{tikzcd}[ampersand replacement=\&,row sep=0.7em,column sep=small]
{} \& {\llap{$\bullet$\hspace{0.8em}}}\arrow[d,mapsto]\\
0\arrow[r] \& 1
\end{tikzcd}}

Je n'ai pas \unsure{vu|eu} d'exemple montrant que la
notion de flèches \uline{loc}. cofiltrantes resp. précofiltrantes
n'est pas stable par changt. de base, et ne voudrais pas m'y attarder, n'ayant pas besoin de
connaître la réponse.

2) Pour que $X$ soit cofiltrante (resp. \ldots),
il f. et s. que $X\longrightarrow e$ le soit. Ainsi la
notion absolue est ramenée à une notion
relative.

3) Si $f:X\to Y$ est cofiltrante, elle est
% archive sheet number: 236
